\documentclass[11pt,a4paper]{article}

% ---------------------------------------------------------
% PAQUETES
% ---------------------------------------------------------
\usepackage[utf8]{inputenc}
\usepackage[T1]{fontenc}
\usepackage[spanish,es-nodecimaldot]{babel}
\usepackage{lmodern}
\usepackage[a4paper,margin=2.5cm]{geometry}
\usepackage{xcolor}
\usepackage{listings}
\usepackage{booktabs}
\usepackage{longtable}
\usepackage{enumitem}
\usepackage{hyperref}
\usepackage{fancyhdr}
\usepackage{tcolorbox}
\usepackage{pmboxdraw}
\usepackage{menukeys}


% ---------------------------------------------------------
% CONFIGURACIÓN DE HYPERREF
% ---------------------------------------------------------
\hypersetup{
    colorlinks=true,
    linkcolor=blue!60!black,
    urlcolor=blue!60!black
}

% ---------------------------------------------------------
% CONFIGURACIÓN DE LISTINGS
% ---------------------------------------------------------
\definecolor{codebg}{RGB}{245,245,245}
\definecolor{codecomment}{RGB}{90,90,90}

\lstdefinestyle{terminal}{
    backgroundcolor=\color{codebg},
    basicstyle=\ttfamily\small,
    frame=single,
    framerule=0.4pt,
    rulecolor=\color{gray!40},
    columns=fullflexible,
    keepspaces=true,
    showstringspaces=false,
    breaklines=true,
    commentstyle=\color{codecomment},
    language=,
    % Replace box-drawing characters with their text equivalents
    extendedchars=true,
    literate=
	    {├}{{\pmboxdrawuni{251C}}}1
	    {└}{{\pmboxdrawuni{2514}}}1
	    {─}{{\pmboxdrawuni{2500}}}1
	    {│}{{\pmboxdrawuni{2502}}}1
	    {┬}{{\pmboxdrawuni{252C}}}1
	    {┤}{{\pmboxdrawuni{2524}}}1
	    {┼}{{\pmboxdrawuni{253C}}}1
	    {á}{{\'a}}1
	    {é}{{\'e}}1
	    {í}{{\'i}}1
	    {ó}{{\'o}}1
	    {ú}{{\'u}}1
	    {Á}{{\'A}}1
	    {É}{{\'E}}1
	    {Í}{{\'I}}1
	    {Ó}{{\'O}}1
	    {Ú}{{\'U}}1
	    {ñ}{{\~n}}1
	    {Ñ}{{\~N}}1
	    {ü}{{\"u}}1
	    {Ü}{{\"U}}1
        {¿}{{\textquestiondown}}1
	    % Ensure LaTeX doesn't misinterpret special characters
}

\lstset{style=terminal}

% ---------------------------------------------------------
% CAJAS
% ---------------------------------------------------------
\newtcolorbox{importante}{
    colback=yellow!10,
    colframe=orange!70!black,
    title=\textbf{Importante},
    fonttitle=\bfseries
}

\newtcolorbox{peligro}{
    colback=red!5,
    colframe=red!70!black,
    title=\textbf{¡Cuidado!},
    fonttitle=\bfseries
}

\newtcolorbox{ejercicio}{
    colback=blue!5,
    colframe=blue!60!black,
    title=\textbf{Ejercicio},
    fonttitle=\bfseries
}

% ---------------------------------------------------------
% ENCABEZADO Y PIE
% ---------------------------------------------------------
\pagestyle{fancy}
\fancyhf{}
\lhead{Terminal de Windows}
\rhead{Guía básica}
\cfoot{\thepage}

% ---------------------------------------------------------
% DOCUMENTO
% ---------------------------------------------------------
\begin{document}

% ---------------------------------------------------------
% PORTADA
% ---------------------------------------------------------
\begin{titlepage}
    \centering

    \vspace*{2cm}

    {\Huge\bfseries Guía básica de la terminal de Windows\par}

    \vspace{1cm}

    {\Large Gestión de directorios y ficheros\par}

    \vspace{1cm}
    
    {\Large Laboratorio de Matemáticas\par}

    \vspace{1cm}

	{Grado en Ingeniería Matemática - Doble Grado en Física e Ingeniería Matemática\par}

    \vspace{1cm}

    {\Large Miguel Beltrá Vidal\par}

    \vfill
    
    \begin{tcolorbox}[width=0.8\textwidth]
    En esta guía aprenderemos a utilizar los principales comandos
    del Símbolo del sistema de Windows para movernos por el sistema
    de archivos y gestionar directorios y ficheros.
    \end{tcolorbox}

    \vfill

    {\large Material práctico para estudiantes\par}

    \vspace{0.5cm}

    {\large \today\par}

\end{titlepage}

\tableofcontents
\newpage

% =========================================================
\section{Introducción}
% =========================================================

La \textbf{terminal} es una herramienta que permite comunicarnos con el sistema operativo mediante comandos de texto. En Windows disponemos principalmente de dos terminales:

\begin{itemize}
    \item El \textbf{Símbolo del sistema} (\texttt{cmd.exe}), que es el
          que utilizaremos en esta guía.
    \item \textbf{PowerShell}, más moderna y potente, que veremos
          brevemente en el apéndice.
\end{itemize}

\noindent Para abrir el Símbolo del sistema podemos pulsar la tecla
\texttt{Windows}, escribir \texttt{cmd} y pulsar \texttt{Intro}.

\noindent En esta guía aprenderemos a:

\begin{itemize}
    \item Saber en qué directorio estamos.
    \item Listar el contenido de un directorio.
    \item Cambiar de directorio y de unidad.
    \item Crear directorios (carpetas).
    \item Crear ficheros.
    \item Copiar ficheros y directorios.
    \item Mover y renombrar ficheros y directorios.
    \item Borrar ficheros y directorios.
    \item Visualizar la estructura de directorios mediante
          \texttt{tree}.
\end{itemize}

\begin{importante}
En Windows los directorios se separan con la barra invertida
(\texttt{\textbackslash}), no con la barra normal, y las rutas comienzan
por la letra de la unidad. Por ejemplo:
\texttt{C:\textbackslash Users\textbackslash alumno}.
\end{importante}

\begin{peligro}
Algunos comandos, especialmente \texttt{del} y \texttt{rmdir /s},
pueden borrar información de forma permanente. Antes de ejecutar un
comando de borrado, comprueba siempre dónde estás y qué vas a borrar.
\end{peligro}

% =========================================================
\section{¿Dónde estoy?}
% =========================================================

Para saber en qué directorio nos encontramos escribimos \texttt{cd}
sin ningún argumento:

\begin{lstlisting}
cd
\end{lstlisting}

\noindent \texttt{cd} significa \textit{change directory}, pero cuando
se utiliza sin argumentos muestra el directorio actual. Por ejemplo:

\begin{lstlisting}
C:\Users\alumno> cd
C:\Users\alumno
\end{lstlisting}

\noindent Esto indica que estamos dentro del directorio
\texttt{C:\textbackslash Users\textbackslash alumno}.

\noindent También podemos utilizar:

\begin{lstlisting}
echo %cd%
\end{lstlisting}

\begin{importante}
Observa que el propio \textit{prompt} del Símbolo del sistema ya nos
muestra el directorio en el que estamos, antes del símbolo
\texttt{>}.
\end{importante}

% =========================================================
\section{Listar el contenido de un directorio}
% =========================================================

Para ver qué hay dentro del directorio actual utilizamos:

\begin{lstlisting}
dir
\end{lstlisting}

\noindent Por ejemplo:

\begin{lstlisting}
C:\Users\alumno> dir
 El volumen de la unidad C no tiene etiqueta.

 Directorio de C:\Users\alumno

17/09/2026  10:12    <DIR>          .
17/09/2026  10:12    <DIR>          ..
17/09/2026  10:12    <DIR>          Descargas
17/09/2026  10:12    <DIR>          Documentos
17/09/2026  10:12    <DIR>          Imagenes
17/09/2026  10:12    <DIR>          Musica
               0 archivos              0 bytes
               6 dirs   35.000.000.000 bytes libres
\end{lstlisting}

\noindent Las carpetas aparecen marcadas con \texttt{<DIR>}.

\subsection{Algunas opciones interesantes}

Para obtener un listado en una sola columna, más compacto:

\begin{lstlisting}
dir /b
\end{lstlisting}

\noindent Para mostrar también los archivos ocultos y del sistema:

\begin{lstlisting}
dir /a
\end{lstlisting}

\noindent Para mostrar solamente los directorios:

\begin{lstlisting}
dir /ad
\end{lstlisting}

\noindent Podemos combinar opciones:

\begin{lstlisting}
dir /a /b
\end{lstlisting}

\begin{ejercicio}
\begin{enumerate}
    \item Abre el Símbolo del sistema.
    \item Comprueba en qué directorio estás con \texttt{cd}.
    \item Lista su contenido con \texttt{dir}.
    \item Utiliza \texttt{dir /b}.
    \item Utiliza \texttt{dir /a}.
\end{enumerate}
\end{ejercicio}

% =========================================================
\section{Cambiar de directorio}
% =========================================================

Para desplazarnos entre directorios utilizamos:

\begin{lstlisting}
cd nombre_carpeta
\end{lstlisting}

\noindent Por ejemplo, si tenemos un directorio llamado
\texttt{Documentos}:

\begin{lstlisting}
cd Documentos
\end{lstlisting}

\noindent Podemos comprobar que hemos cambiado de directorio:

\begin{lstlisting}
cd
\end{lstlisting}

\subsection{Volver al directorio anterior}

Para subir un nivel en la estructura de directorios:

\begin{lstlisting}
cd ..
\end{lstlisting}

\noindent Por ejemplo, si estamos en \verb|C:\Users\alumno\Documentos|, después de ejecutar:

\begin{lstlisting}
cd ..
\end{lstlisting}

\noindent estaremos en: \verb|C:\Users\alumno|.


\subsection{Ir al directorio personal}

El directorio personal del usuario está guardado en la variable
\texttt{\%USERPROFILE\%}:

\begin{lstlisting}
cd %USERPROFILE%
\end{lstlisting}

\noindent Por ejemplo, esto nos llevaría a
\texttt{C:\textbackslash Users\textbackslash alumno}.

\subsection{Ir al directorio raíz}

El directorio raíz de la unidad actual se representa mediante
\texttt{\textbackslash}:

\begin{lstlisting}
cd \
\end{lstlisting}

\noindent Por ejemplo, si estamos en \verb|C:\Users\alumno|, y ejecutamos
\begin{lstlisting}
cd \
\end{lstlisting}
cambiaremos a \verb|C:\|.

\subsection{Cambiar de unidad}

En Windows cada unidad (disco duro, pendrive, etc.) tiene su propia
letra. Para cambiar de unidad escribimos la letra seguida de dos
puntos:

\begin{lstlisting}
D:
\end{lstlisting}

\noindent Si queremos cambiar de unidad y de directorio a la vez,
necesitamos la opción \texttt{/d}:

\begin{lstlisting}
cd /d D:\practicas
\end{lstlisting}

\begin{peligro}
Sin la opción \texttt{/d}, el comando \texttt{cd} no cambia de unidad:
cambiará el directorio actual de esa otra unidad, pero seguiremos
trabajando en la unidad en la que estábamos.
\end{peligro}

\begin{ejercicio}
Partiendo de tu directorio personal:

\begin{enumerate}
    \item Entra en \texttt{Documentos}.
    \item Comprueba dónde estás con \texttt{cd}.
    \item Vuelve al directorio anterior con \texttt{cd ..}.
    \item Entra de nuevo en \texttt{Documentos}.
    \item Vuelve a tu directorio personal utilizando \verb|cd %USERPROFILE%|.
\end{enumerate}
\end{ejercicio}

% =========================================================
\section{Crear directorios}
% =========================================================

Para crear una directorio (carpeta) utilizamos:

\begin{lstlisting}
mkdir nombre_carpeta
\end{lstlisting}

\noindent \texttt{mkdir} significa \textit{make directory}. También
podemos escribir su forma abreviada \texttt{md}. Por ejemplo, si
queremos crear el directorio \texttt{practica}:

\begin{lstlisting}
mkdir practica
\end{lstlisting}

\noindent Podemos comprobar que se ha creado:

\begin{lstlisting}
dir
\end{lstlisting}

\subsection{Crear varias carpetas}

También podemos crear varias carpetas en un único comando:

\begin{lstlisting}
mkdir documentos imagenes programas
\end{lstlisting}

\subsection{Crear una estructura completa}

En Windows, \texttt{mkdir} crea automáticamente todos los directorios
intermedios que hagan falta:

\begin{lstlisting}
mkdir windows\documentos\apuntes
\end{lstlisting}

\begin{ejercicio}
Dentro de tu directorio personal crea la siguiente estructura:

\begin{lstlisting}
windows\
├── documentos\
├── imagenes\
├── programas\
└── ejercicios\
\end{lstlisting}

Puedes hacerlo mediante:

\begin{lstlisting}
mkdir windows
cd windows
mkdir documentos imagenes programas ejercicios
\end{lstlisting}
\end{ejercicio}

% =========================================================
\section{Crear ficheros}
% =========================================================

Una forma sencilla de crear un fichero vacío con el nombre
\texttt{nombre\_fichero} es mediante:

\begin{lstlisting}
type nul > nombre_fichero
\end{lstlisting}

\noindent Por ejemplo:

\begin{lstlisting}
type nul > notas.txt
\end{lstlisting}

\noindent Comprobamos que existe:

\begin{lstlisting}
dir
\end{lstlisting}

\noindent Si queremos crear el fichero con algún contenido, podemos
utilizar \texttt{echo}:

\begin{lstlisting}
echo Mis apuntes de Windows > notas.txt
\end{lstlisting}

\begin{importante}
El símbolo \texttt{>} redirige la salida de un comando hacia un
fichero. Si el fichero ya existía, \textbf{su contenido se
sobrescribe}. Para añadir texto al final sin borrar lo anterior se
utiliza \texttt{>}\texttt{>}.
\end{importante}

\noindent Para crear varios ficheros necesitamos un comando por
fichero:

\begin{lstlisting}
type nul > uno.txt
type nul > dos.txt
type nul > tres.txt
\end{lstlisting}

\begin{ejercicio}
Dentro del directorio \texttt{documentos}, crea estos ficheros:

\begin{lstlisting}
apuntes.txt
tareas.txt
examen.txt
\end{lstlisting}

Comprueba que se han creado utilizando:

\begin{lstlisting}
dir
\end{lstlisting}
\end{ejercicio}

% =========================================================
\section{Copiar ficheros}
% =========================================================

Para copiar un fichero utilizamos:

\begin{lstlisting}
copy origen destino
\end{lstlisting}

\noindent Por ejemplo:

\begin{lstlisting}
copy apuntes.txt copia_apuntes.txt
\end{lstlisting}

\noindent Ahora tendremos:

\begin{lstlisting}
apuntes.txt
copia_apuntes.txt
\end{lstlisting}

\noindent También podemos copiar un fichero a otro directorio:

\begin{lstlisting}
copy apuntes.txt ..\
\end{lstlisting}

\noindent En este caso, \texttt{..\textbackslash} representa el
directorio padre.

\subsection{Copiar un directorio}

El comando \texttt{copy} no puede copiar directorios. Para ello, utilizamos \texttt{xcopy}:

\begin{lstlisting}
xcopy documentos documentos_copia /e /i
\end{lstlisting}

\noindent Aquí:

\begin{itemize}
    \item \texttt{/e} copia el directorio de forma
          \textbf{recursiva}, incluyendo subdirectorios (también los
          vacíos).
    \item \texttt{/i} indica que el destino es un directorio, aunque
          todavía no exista.
\end{itemize}

\noindent Una alternativa más moderna y potente es \texttt{robocopy}:

\begin{lstlisting}
robocopy documentos documentos_copia /e
\end{lstlisting}

% =========================================================
\section{Mover y renombrar}
% =========================================================

Para mover un fichero utilizamos:

\begin{lstlisting}
move origen destino
\end{lstlisting}

\noindent Por ejemplo, si queremos mover el fichero \texttt{apuntes.txt} del directorio actual al directorio padre, emplearemos el comando:
\begin{lstlisting}
move apuntes.txt ..\
\end{lstlisting}



\subsection{Mover un fichero a otra carpeta}

Por ejemplo:

\begin{lstlisting}
move examen.txt ..\examenes\
\end{lstlisting}

\subsection{Mover un directorio}

El comando \texttt{move} también sirve para mover directorios
completos, sin necesidad de ninguna opción adicional. Por ejemplo, si queremos mover el directorio \texttt{imagenes} dentro del directorio \texttt{documentos}, utilizaremos el comando:

\begin{lstlisting}
move imagenes documentos\
\end{lstlisting}

\subsection{Renombrar un fichero}

Para cambiar el nombre utilizamos \texttt{ren} (o \texttt{rename}):

\begin{lstlisting}
ren tareas.txt tareas_windows.txt
\end{lstlisting}

\noindent No estamos creando un fichero nuevo: simplemente estamos
cambiando su nombre.

\subsection{Renombrar un directorio}

También podemos hacer:

\begin{lstlisting}
ren imagenes fotos
\end{lstlisting}

\begin{importante}
A diferencia de Linux, en Windows mover y renombrar son dos comandos
distintos: \texttt{move} para mover y \texttt{ren} para renombrar.
Aun así, \texttt{move} también puede renombrar si el destino está en
el mismo directorio.
\end{importante}

% =========================================================
\section{Borrar ficheros}
% =========================================================

Para borrar un fichero utilizamos:

\begin{lstlisting}
del nombre_fichero
\end{lstlisting}

\noindent Por ejemplo:

\begin{lstlisting}
del tres.txt
\end{lstlisting}

\noindent El fichero se eliminará. También podemos utilizar
\texttt{erase}, que hace exactamente lo mismo.

\begin{peligro}
\texttt{del} no envía el fichero a la Papelera de reciclaje. La
eliminación es permanente.
\end{peligro}

\noindent Podemos pedir confirmación antes de borrar:

\begin{lstlisting}
del /p tres.txt
\end{lstlisting}

\noindent Para borrar varios ficheros a la vez podemos utilizar
comodines. Por ejemplo, todos los ficheros con extensión
\texttt{.txt}:

\begin{lstlisting}
del *.txt
\end{lstlisting}

\begin{peligro}
El comodín \texttt{*} puede afectar a muchos más ficheros de los que
esperamos. Antes de borrar con comodines, comprueba primero qué
ficheros coinciden utilizando \texttt{dir}. Por ejemplo:
\texttt{dir *.txt}.
\end{peligro}

% =========================================================
\section{Borrar directorios}
% =========================================================

Para borrar un directorio vacío podemos utilizar:

\begin{lstlisting}
rmdir nombre_directorio
\end{lstlisting}

\noindent Por ejemplo:

\begin{lstlisting}
rmdir carpeta_vacia
\end{lstlisting}

\noindent \texttt{rmdir} (o su forma abreviada \texttt{rd})
solamente puede eliminar directorios vacíos.

\subsection{Borrar un directorio y todo su contenido}

Para eliminar un directorio que contiene ficheros u otros
directorios utilizamos la opción \texttt{/s}:

\begin{lstlisting}
rmdir /s nombre_directorio
\end{lstlisting}

\noindent Por ejemplo, si ejecutamos:

\begin{lstlisting}
rmdir /s ejercicios
\end{lstlisting}

\noindent se eliminará el directorio \texttt{ejercicios} y su contenido. Windows nos pedirá confirmación:

\begin{lstlisting}
ejercicios, ¿Está seguro (S/N)?
\end{lstlisting}

\noindent Si añadimos la opción \texttt{/q} (\textit{quiet}), no se
pedirá confirmación:

\begin{lstlisting}
rmdir /s /q ejercicios
\end{lstlisting}

\begin{peligro}
\texttt{rmdir /s /q} puede eliminar muchos archivos de una sola vez y
sin avisar. Antes de utilizarlo, comprueba siempre el contenido y
asegúrate de que estás en el directorio correcto.
\end{peligro}

% =========================================================
\section{Visualizar la estructura con \texttt{tree}}
% =========================================================

El comando \texttt{tree} permite visualizar los directorios formando
una estructura de árbol. Por ejemplo, tras escribir:

\begin{lstlisting}
tree
\end{lstlisting}

\noindent Podríamos obtener:

\begin{lstlisting}
Listado de rutas de carpetas
C:.
├───documentos
├───ejercicios
├───imagenes
└───programas
\end{lstlisting}

\subsection{Mostrar también los ficheros}

Por defecto \texttt{tree} muestra solamente los directorios. Para que
muestre también los ficheros utilizamos la opción \texttt{/f}:

\begin{lstlisting}
tree /f
\end{lstlisting}

\noindent Por ejemplo:

\begin{lstlisting}
C:.
├───documentos
│       apuntes.txt
│       examen.txt
│       tareas.txt
│
├───ejercicios
├───imagenes
└───programas
\end{lstlisting}

\subsection{Utilizar caracteres ASCII}

Si los caracteres del árbol no se muestran correctamente, podemos
utilizar la opción \texttt{/a}, que dibuja el árbol con caracteres
ASCII sencillos:

\begin{lstlisting}
tree /a
\end{lstlisting}

\noindent Por ejemplo:

\begin{lstlisting}
C:.
+---documentos
+---ejercicios
+---imagenes
\---programas
\end{lstlisting}

\begin{importante}
A diferencia de Linux, \texttt{tree} viene siempre incluido en
Windows, por lo que no hace falta instalar nada. Eso sí, ten en
cuenta que \texttt{tree} muestra solamente directorios y
\texttt{tree /f} muestra directorios y ficheros.
\end{importante}




\section{Limpiar la terminal}

Si queremos limpiar la terminal, podemos utilizar el comando:

\begin{lstlisting}
cls
\end{lstlisting}




\section{Autocompletado}

La mejor forma de escribir el nombre de un directorio o fichero es dejando que sea la propia terminal la que lo escriba. La terminal tiene la función de autocompletado con la tecla \texttt{tab} \keys{\tab}. Por ejemplo, supongamos que nuestro directorio actual tiene la forma:

\begin{lstlisting}
windows\
├── documentos\
├── imagenes\
├── programas\
└── ejercicios\
\end{lstlisting}

Si escribimos 

\begin{lstlisting}
cd do
\end{lstlisting}

\noindent y presionamos la tecla \texttt{tab} \keys{\tab}, entonces se autocompletará como

\begin{lstlisting}
cd documentos
\end{lstlisting}


\newpage
\appendix
% =========================================================
\section{Resumen de comandos}
% =========================================================

\begin{longtable}{p{0.20\textwidth}p{0.36\textwidth}p{0.34\textwidth}}
\toprule
\textbf{Comando} & \textbf{Función} & \textbf{Ejemplo} \\
\midrule
\endhead

\texttt{cd} &
Muestra el directorio actual &
\texttt{cd} \\

\texttt{dir} &
Lista el contenido &
\texttt{dir} \\

\texttt{dir /b} &
Lista en una sola columna &
\texttt{dir /b} \\

\texttt{dir /a} &
Incluye archivos ocultos &
\texttt{dir /a} \\

\texttt{cd} &
Cambia de directorio &
\texttt{cd documentos} \\

\texttt{cd ..} &
Sube al directorio padre &
\texttt{cd ..} \\

\texttt{cd \textbackslash} &
Va a la raíz de la unidad &
\texttt{cd \textbackslash} \\

\texttt{cd /d} &
Cambia de unidad y directorio &
\texttt{cd /d D:\textbackslash datos} \\

\texttt{mkdir} &
Crea un directorio &
\texttt{mkdir fotos} \\

\texttt{type nul >} &
Crea un fichero vacío &
\texttt{type nul > notas.txt} \\

\texttt{copy} &
Copia un fichero &
\texttt{copy a.txt b.txt} \\

\texttt{xcopy /e /i} &
Copia un directorio &
\texttt{xcopy fotos fotos2 /e /i} \\

\texttt{move} &
Mueve un fichero/directorio &
\texttt{move a.txt documentos\textbackslash} \\

\texttt{ren} &
Renombra &
\texttt{ren a.txt notas.txt} \\

\texttt{del} &
Borra un fichero &
\texttt{del notas.txt} \\

\texttt{rmdir} &
Borra un directorio vacío &
\texttt{rmdir temporal} \\

\texttt{rmdir /s} &
Borra un directorio y su contenido &
\texttt{rmdir /s temporal} \\

\texttt{tree} &
Muestra los directorios en árbol &
\texttt{tree} \\

\texttt{tree /f} &
Incluye también los ficheros &
\texttt{tree /f} \\

\texttt{cls} &
Limpia la pantalla &
\texttt{cls} \\

\bottomrule
\end{longtable}

% =========================================================






\newpage
\section{Equivalencias entre Linux, CMD y PowerShell}
% =========================================================

PowerShell es la otra terminal disponible en Windows. Para abrirla, pulsamos la tecla \texttt{Windows}, escribimos \texttt{powershell} y pulsamos \texttt{Intro}.

Una de sus ventajas es que incorpora \textit{alias} con los nombres habituales de Linux, por lo que muchos comandos de esta guía funcionan igual que allí.

\begin{longtable}{p{0.2\textwidth}p{0.32\textwidth}p{0.40\textwidth}}
\toprule
\textbf{Linux} & \textbf{CMD} & \textbf{PowerShell} \\
\midrule
\endhead

\texttt{pwd} &
\texttt{cd} &
\texttt{pwd} o \texttt{Get-Location} \\

\texttt{ls} &
\texttt{dir} &
\texttt{ls} o \texttt{Get-ChildItem} \\

\texttt{ls -la} &
\texttt{dir /a} &
\texttt{ls -Force} \\

\texttt{cd} &
\texttt{cd} &
\texttt{cd} o \texttt{Set-Location} \\

\texttt{cd \textasciitilde} &
\texttt{cd\,\,\%USERPROFILE\%} &
\texttt{cd \textasciitilde} \\

\texttt{mkdir} &
\texttt{mkdir} &
\texttt{mkdir} o \texttt{New-Item -Type Directory} \\

\texttt{touch} &
\texttt{type nul >} &
\texttt{New-Item notas.txt} \\

\texttt{cp} &
\texttt{copy} &
\texttt{cp} o \texttt{Copy-Item} \\

\texttt{cp -r} &
\texttt{xcopy /e /i} &
\texttt{cp -Recurse} \\

\texttt{mv} &
\texttt{move} / \texttt{ren} &
\texttt{mv} o \texttt{Move-Item} \\

\texttt{rm} &
\texttt{del} &
\texttt{rm} o \texttt{Remove-Item} \\

\texttt{rm -r} &
\texttt{rmdir /s} &
\texttt{rm -Recurse} \\

\texttt{clear} &
\texttt{cls} &
\texttt{clear} o \texttt{cls} \\

\bottomrule
\end{longtable}

\begin{importante}
En PowerShell se pueden utilizar tanto la barra invertida
(\texttt{\textbackslash}) como la barra normal (\texttt{/}) para
separar directorios, lo que facilita el paso de un sistema a otro.
\end{importante}

\end{document}